% Translated from: thesis/chapter_01_绪论.md
% Target journal: General
% Generated: 2026-04-03
\chapter{Introduction}
\setcounter{figure}{0}
\setcounter{equation}{0}
\renewcommand{\thefigure}{\thechapter-\arabic{figure}}
\renewcommand{\theequation}{\thechapter-\arabic{equation}}

\section{Research Background and Significance}

General Game Playing (GGP) aims to enable an agent to receive formal rule descriptions at runtime and perform game reasoning, action selection, and strategy generation without relying on domain-specific priors. Unlike traditional game programs that are designed around handcrafted features, evaluation functions, and search rules for a single board game, GGP emphasizes unification across the rule layer, state-machine layer, and decision layer. Its value lies not only in game intelligence itself, but also in validating general reasoning, general search, and cross-task learning capabilities \cite{genesereth2005general,genesereth2026general}. Therefore, GGP has long been regarded as an important testbed for studying how agents can rapidly form effective decision-making ability in unfamiliar rule environments.

In existing GGP research, Monte Carlo Tree Search (MCTS) is widely adopted because it does not depend on strong domain knowledge and can perform direct forward simulation on a given state machine \cite{browne2012survey}. Especially in settings where rules are unknown in advance but executable, MCTS incrementally builds a search tree through selection, expansion, simulation, and backpropagation, providing a relatively robust decision foundation for agents. Furthermore, enhancement techniques such as heuristic rollouts, historical statistics, and neural-network-based evaluation have continuously been introduced into MCTS to improve search efficiency and decision quality in complex game environments \cite{finnsson2008simulation,anthony2017thinking,silver2018general}. At the same time, considering a key property of GGP, namely that games are specified at runtime and are difficult to tune precisely offline, Sironi et al. proposed an adaptive MCTS approach that adjusts search-control parameters online for a specific game. This also indicates the practical complexity of search enhancement under unified rule environments \cite{sironi2020selfadaptive}.

However, search enhancement in GGP is not as straightforward as in fixed game domains. First, substantial differences in rule structure, action space, and state representation across games make it difficult to transfer heuristics that work well for a single game. Second, as research moves from perfect-information board games to more complex rule settings involving hidden information, knowledge constraints, and even communication behavior, the difficulty of unified representation and unified reasoning further increases \cite{partridge2022hidden,zammit2023communication}. Third, pure search methods are often constrained by simulation overhead under limited decision time, and the search budget is further compressed when rule interpretation and state transition themselves already consume considerable computation. Finally, although neural value methods create new opportunities for search enhancement, they still face difficulties in GGP, including unified representation, coarse supervision design, insufficient cross-game generalization, and unstable integration with search \cite{goldwaser2020deep,mcewan2022knowledge,maras2024fast}. These points show that, in a unified rule environment, building a value-evaluation module that can both adapt to multiple games and genuinely improve search performance under limited budgets remains an important research problem.

From an engineering and application perspective, this problem has dual significance. On one hand, if a reusable value-correction mechanism can be realized within a unified GGP framework, reliance on handcrafted knowledge for specific games may be reduced, thereby improving adaptation to new rule environments. On the other hand, if a value model can be translated into actual search gains under limited time budgets, it supports the practical feasibility of the technical route of rule-driven execution, search-distillation supervision, and residual value correction. Recent discussions on search, planning, learning, and interpretability have also continued to list GGP-related topics as important directions, indicating sustained research interest and room for extension in this area \cite{soemers2024gametable}. Therefore, search enhancement research around GGP and MCTS is important not only theoretically, but also for system implementation and experimental validation.

\section{Research Questions and Challenges}

Based on the background above, this thesis focuses on the following question: under a unified GGP framework, how can search distillation and residual learning be used to build a neural value model that can be stably integrated into MCTS, so that the quality of leaf-node evaluation and real match performance can be improved under limited search budgets without breaking the logic of rule-driven execution.

This question involves at least three levels of challenges. First, the core premise of GGP is that rules are given at runtime, so the input representation of the value model must remain as general as possible and cannot rely on game-specific feature engineering; otherwise, cross-game reuse is difficult. Second, if only terminal win/loss outcomes or weak-search returns are used as supervision, value learning often suffers from high supervisory noise and mismatch with actual search demands, making it difficult for the model to learn evaluation signals that genuinely help search. Third, even after a value model is trained, stable integration into MCTS remains nontrivial: what should neural evaluation replace, what should it complement, and how can its real contribution be validated under a unified search skeleton.

As indicated by the existing chapter content, this work does not attempt to directly train a general game policy network that operates independently of search. Instead, it narrows the problem to value correction, which is more aligned with current project conditions and research goals. In other words, rather than fully rebuilding the entire decision stack, this thesis starts from low-cost search estimates, uses stronger search to generate teacher signals, and then learns correction terms so that final predictions are closer to strong-search evaluations and can be converted into online playing strength. This formulation preserves the generality requirements of GGP while remaining consistent with the system implementation and experimental design in later chapters.

\section{Research Approach and Main Contributions}

To address the above problem, this thesis proposes the Search-Distilled Residual Progress-guided Value (SDRPV) method and builds a complete experimental loop from rule execution and sample generation to model training and online evaluation. The overall idea is as follows: first, run self-play and search over a unified state machine to construct both low-cost baseline values and high-budget teacher values for intermediate states; second, use a unified board-token state representation and a Transformer value network to learn corrections to baseline estimates; finally, integrate the trained value model into a shared MCTS kernel and evaluate its practical effect against pure search, heuristic search, and neural-value baselines under a fixed-time budget.

The main research content and contributions of this thesis can be summarized in four aspects.

\begin{enumerate}
    \item It establishes unified rule execution and state-machine foundations for GGP. Using GDL/KIF rule descriptions as input, this work implements a unified state-machine interface through rule parsing and logical reasoning pipelines, so that \texttt{breakthrough}, \texttt{connectFour}, and \texttt{hex} can all perform legal-action queries, state transitions, and terminal scoring under a consistent environment semantics. This provides a unified execution basis for subsequent method comparison.

    \item It proposes the SDRPV method centered on search distillation and residual correction. Instead of directly regressing the complete value, this work constructs both a low-cost baseline value and a high-budget teacher value for the same state and uses their difference as the core learning target, turning value learning into the more stable task of correcting existing estimation bias. This design makes supervision signals closer to real search demands and aligns the model objective with leaf-evaluation improvement under limited budgets.

    \item It designs and implements a unified scheme for state encoding, value estimation, and search integration. This work adopts a board-token-based general representation and outputs scalar value estimates with a Transformer architecture, allowing the neural evaluator to be integrated into a shared MCTS kernel while preserving cross-game consistency. Compared with directly replacing the rule system, this approach consistently preserves the dominant role of rule-grounded truth in terminal judgment, thereby clarifying the engineering boundary between rule execution and neural evaluation.
    
    \item It completes systematic experimental validation. This work uses \texttt{connectFour} self-play samples to build training data and performs fixed-time online evaluation on \texttt{breakthrough}, \texttt{connectFour}, and \texttt{hex}. Results show that SDRPV outperforms baseline value \texttt{b} and the teacher-only route on offline metrics. In the online baseline benchmark, SDRPV achieves an overall win rate of \texttt{0.7133} against both \texttt{pure_mct} and \texttt{heuristic_mcts} across the three games, and also obtains an overall positive edge of \texttt{0.5467} against \texttt{baseline_token_transformer}. The ablation benchmark further indicates that teacher supervision and the residual objective jointly contribute to final performance gains. These results demonstrate that the proposed technical route can be converted into real search gains within a unified GGP framework.
\end{enumerate}

Overall, the focus of this thesis is not to propose a brand-new game-playing agent that is fully detached from existing search frameworks. Instead, it targets value correction under a unified rule environment and develops a methodological route with clear supervision sources, explicit system boundaries, and verifiable experimental gains. This is also why, compared with general neural game-playing studies, this thesis places stronger emphasis on search enhancement in general-rule settings.
